\documentclass[../Main.tex]{subfiles}

\begin{document}

\section{Evaluation} \label{sec:Evaluation}

Das Ziel des Playtests ist, die acht Designvorschläge (V1--V8), die wir aus unseren Spieleanalysen und der Leveldesign-Recherche hergeleitet und in Kapitel~\ref{sec:Umsetzung} umgesetzt haben, zu überprüfen, also zu prüfen, ob diese die beabsichtigten Wirkungen auf Spielfluss, Schwierigkeit und Spielgefühl tatsächlich erzeugen.

Wir sind für den playtest wie in kapitel 2.6 beschrieben vorgegenagen. also Pre-Questionnaire zum einschätzen der testperson dann der playtest der aufgenommen wurde, und dann das Postquestionnaire und das Semistrukturierte interview für tiefer gehende fragen falls vorhanden.

Das folgende kapitel beschreibt den ablauf des Playtests und geht auf für die Erkenntnisfrage relevante Ergebnisse ein. Darauf folgt eine diskussion bei der wir uns an jedem unserer Designvorschläge abarbeiten und unsere Empfehlung für weitere Playtests abgeben.

\subsection{Durchführung und Testpersonen} \label{sec:Durchfuehrung}
Der playtest wurde immer von Rodin Arslan durchgefüht, Zhile Zhu war auch hin und wieder dabei. Die playtest wurden immer auf maus und tastatur durchgespielt. 3 playtest wurden via Discord gespielt mit screenshare und 2 in person. die playtests wurden alle auf der gleichen export version gespielt. Mac oder Windows. Die spieler haben immer ihre eigene und gewohnte hardware benutzt. Es wurde währenddessen kein smalltalk gefürt. allerdings haben wir mit den testpersonen in zeiten von entspannung über das spiel geredet. wenn spieler eine frage hatten zum combat wurde sie beantwortet oder wenn sie was vergessen hatten. Allerdings wurde ihnen ganz am anfang von uns gezeigt, wie man was machen kann und was es alles gibt.
\begin{itemize}[leftmargin=1.5cm]
    \item \textcolor{red}{\textbf{Die folgenden vier Absätze wiederholen den geplanten Ablauf aus Abschnitt~\ref{sec:VorgehenPlaytest} und sind umgangssprachlich.} Ersetzen durch die \emph{tatsächliche} Durchführung: Datum/Ort, Hardware und Eingabe, Build-Version, Dauer, wer beobachtet hat, Abweichungen vom Plan (mündliche Erklärung, weil es kein Tutorial gab; nur spontane Reaktionen statt durchgehendem Think-aloud).}
    \item \textcolor{red}{Anonymisierte Tabelle der Testpersonen TP1--TP5: Spielzeit pro Woche (P1), Lieblingsgenres (P2), Selbsteinschätzung (P5), Spielstil (P6).}
\end{itemize}

Zunächst wird dem Playtester das Pre-Questionnaire ausgehändigt und gebeten, dass es ausgefüllt wird. Bei Verständnisfragen kann man antworten.

Danach kommt der eigentliche Playtest. Unser spiel eignet sich nicht für dien üblichen think aloud prozess, da dieser das verhalten des spielers ändern/beeinträchtigen kann. Deshalb haben wir uns überlegt, dass sie instruirt werden so mit ausrufen zu arbeiten, wenn sie überrascht sind oder so. Alles wird aufgenommen

Danach kommt das Post-Questionnaire. Wie gehabt beim Pre.

Zum schluss gibt es noch das Semistrukturierte Interview, welches aufschluss über schlüsselmomente liefern kann. Dabei wird die aufnahme abgespielt.

\subsection{Ergebnisse}
Die einzelnen antworten der testpersonen sind im anhang. Hier sind die durchschnittswerte und antworten.

\subsubsection{Quantitative Ergebnisse} \label{sec:ErgQuant}

\begin{itemize}[leftmargin=1.5cm]
    \item \textcolor{red}{Tabelle Q1--Q11 $\times$ TP mit Mittelwert, Median und Spannweite; ein Diagramm. Solange TP1 fehlt: \enquote{n = 4, TP1 folgt}. Keine festen Zahlen im Fliesstext, auf die Tabelle verweisen.}
\end{itemize}

\subsubsection{Qualitative Ergebnisse} \label{sec:ErgQual}

\begin{itemize}[leftmargin=1.5cm]
    \item \textcolor{red}{Pro Designvorschlag V1--V8: was Beobachtung (B), offene Antworten (Q12--Q14) und Interview (I) zeigen, mit Verweis auf den Anhang (TP/Item).}
    \item \textcolor{red}{Danach allgemeine Auffälligkeiten (B10--B14), z.\,B. Türen, fehlendes Tutorial, Steuerung.}
\end{itemize}

\subsection{Diskussion}

\subsubsection{Interpretation pro Designvorschlag} \label{sec:Interpretation}

\begin{itemize}[leftmargin=1.5cm]
    \item \textcolor{red}{Pro V1--V8: bestätigt / teilweise bestätigt / nicht bestätigt, mit Begründung aus den Daten und Bezug zur Recherche. Abweichungen zwischen Vorschlag und Umsetzung (z.\,B. V2, V7) hier diskutieren. Selbst schreiben.}
\end{itemize}

\subsubsection{Gesamteindruck: Spielgefühl, Spielfluss und Schwierigkeit} \label{sec:Gesamteindruck}

\begin{itemize}[leftmargin=1.5cm]
    \item \textcolor{red}{Auswertung von Q9--Q11 und B14 zusammen mit den Ergebnissen pro Vorschlag. Das ist die Brücke zur Erkenntnisfrage in Kapitel~\ref{sec:Schlussfolgerungen}.}
\end{itemize}

\subsection{Verbesserungsvorschläge} \label{sec:Verbesserungen}

\begin{itemize}[leftmargin=1.5cm]
    \item \textcolor{red}{Konkrete Verbesserungen, jede mit Beleg aus den Daten (z.\,B. akustische Angriffsankündigung, Tutorial, Signal beim Öffnen der Tür, abbrechbare Combos, Block). Laut Vertrag müssen sie nicht mehr umgesetzt werden.}
\end{itemize}

\end{document}
